% ═══════════════════════════════════════════════════════════════════
\subsection{Validation Protocol}
\label{sec:validation}
% ═══════════════════════════════════════════════════════════════════

The reference paper proposes to rank candidate docking sites with a SWAP test over the
second-encoding probability amplitudes. A SWAP test, however, returns a similarity in
\emph{encoding space}: it can tell that two windows encode similarly, but not that a
window is spatially correct. Establishing whether the pipeline identifies the binding
site therefore requires an external, geometric ground truth, which we take to be the
crystallographic pose of the co-crystallised ligand. This section defines the protocol;
Section~\ref{sec:results} reports its outcome.

\paragraph{Benchmark set.}
Evaluation runs in benchmark mode (\texttt{--ligand-pdb}) on the $29$ pockets of the
$31$-pocket set for which the matching full structure is available (\texttt{1a08} and
\texttt{1ddm} lack one). The targets are drawn from the Astex Diverse Set plus
hand-picked complexes, and span diverse protein families and binding-site geometries:
\textbf{1GM8, 1GPK, 1HNN, 1HQ2, 1HSG, 1IG3, 1K3U, 1L2S, 1MMV, 1N2V, 1N46, 1OPK, 1OYT,
1Q1G, 1R58, 1R9O, 1S19, 1STP, 1T46, 1TZ8, 1U1C, 1XM6, 1Y6B, 1YGC, 1YV3, 1YWR, 2BSM,
3PTB} and \textbf{4DFR}.

\paragraph{Distance metric.}
For a matching window of $k$ sites, the \emph{candidate centroid} is the arithmetic mean
of the 3D coordinates of its constituent pharmacophoric sites, which are carried
unchanged from the classical extraction stage. The \emph{true ligand centroid} is
computed from the experimental \texttt{HETATM} coordinates of the full structure. Because
asymmetric units frequently contain several bound copies of the same ligand, a centroid
is computed for every copy and the reported metric is the \emph{minimum} Euclidean
distance to any valid copy. This avoids penalising a prediction that correctly targets a
secondary but genuine binding pocket.

\paragraph{Percentile normalisation.}
Raw distances in \AA{} are not comparable across pockets of different size: an $8$~\AA{}
error is near-random in a small pocket and good in a large one. Every candidate is
therefore also reported as its \emph{percentile rank} within the distribution of the
distances of \emph{all} sliding windows of the same pocket, where $0\%$ is the closest
window the representation makes available at all, and $50\%$ is chance. The percentile
is the primary metric throughout Section~\ref{sec:results}; absolute distances are
reported alongside it for interpretability.

\paragraph{Random baseline.}
A percentile is not by itself a sufficient reference, because the pipeline returns
several candidates per pocket and reporting the best of them is a minimum-of-$n$
statistic, which is biased low even for a search that carries no information. For $n$
windows drawn uniformly at random, the expected percentile of the best is
\begin{equation}
    \mathbb{E}\!\left[\min\nolimits_n\right] = \frac{1}{n+1}.
    \label{eq:baseline}
\end{equation}
The baseline is evaluated per pocket with that pocket's own candidate count $n$ and then
averaged, so that the comparison is like-for-like. The ``best of candidates'' row of
Table~\ref{tab:percentile} must be read against this baseline; the ``top-1'' row, which
selects a single window before seeing the ground truth, is read against $50\%$.

\paragraph{Two matching regimes.}
Both variants of the search are evaluated under the identical protocol:
\begin{enumerate}[leftmargin=1.6em]
    \item the \textbf{exact} regime, in which a window matches when its first-encoding
          bitstring equals the ligand's, as prescribed by the reference paper;
    \item the \textbf{amplitude} regime, in which exact equality is replaced by a
          Euclidean distance in the second-encoding amplitude space --- the classical
          counterpart of the paper's SWAP-test similarity.
\end{enumerate}

\paragraph{Classical measurement of a quantum search.}
Both regimes are measured classically rather than by executing the circuit on every
pocket. This is exact, not an approximation. The oracle
$\hat{O} = I - 2\ket{x_{\text{lig}}}\!\bra{x_{\text{lig}}}$ implements equality against a
fixed bitstring, so the set of windows that satisfy it is fully determined by the
encoding \emph{before} any circuit is prepared; amplitude amplification changes which
element of that set is measured, never the set itself. The classical enumeration
therefore characterises the ceiling of the method exactly, while making it feasible to
sweep window sizes and matching regimes over the whole benchmark --- which the dense
unitary synthesis of Section~\ref{sec:grover-search} would not allow.


% ═══════════════════════════════════════════════════════════════════
\section{Results}
\label{sec:results}
% ═══════════════════════════════════════════════════════════════════

This section reports the measurements only; their interpretation is deferred to
Section~\ref{sec:why}.

% ───────────────────────────────────────────────────────────────────
\subsection{The search does not localise the binding site}
\label{sec:res-accuracy}

At the default configuration ($k=3$ sites, $6$ qubits), $23$ of the $29$ pockets produce
candidates. The $6$ that do not all fail for the same reason: no window clears the
uniform $1/N$ probability threshold. Pockets that do produce candidates produce few:
$2.43$ on average, never more than $3$.

The median distance from the top-ranked candidate to the true ligand centroid is
$7.75$~\AA{}, and only $2$ of the $23$ cases fall within a $5.0$~\AA{} threshold
(a third, 1OYT, lands on it at $5.0005$~\AA{})
(Table~\ref{tab:distances}). Table~\ref{tab:percentile} normalises these distances
against the pocket-specific random baseline of Equation~\eqref{eq:baseline}.

\begin{table}[htbp]
\centering
\caption{Percentile ranking of the Grover candidate windows relative to all possible
windows of the same pocket ($k=3$). Lower is better; $50\%$ is chance.}
\label{tab:percentile}
\begin{tabular}{lccc}
\toprule
\textbf{Metric} & \textbf{Mean percentile} & \textbf{Median} & \textbf{Random baseline} \\
\midrule
Top-1 by interactivity score   & $47.2\%$ & $42.0\%$ & $50.0\%$ \\
Best of the candidates returned & $32.2\%$ & $26.2\%$ & $30.8\%$ \\
\bottomrule
\end{tabular}
\end{table}

Two observations follow directly from the table. First, the candidate \emph{set} is
statistically indistinguishable from drawing the same number of windows at random:
$32.2\%$ against a $30.8\%$ baseline is well within the spread of the benchmark, and on
the wrong side of it. Second, the top-1 candidate is \emph{worse} than the best of the
candidates returned ($47.2\%$ vs.\ $32.2\%$), which means that ranking by aggregate
$h/hb$ interactivity does not merely fail to help --- it actively degrades the selection.
The effect is large on individual targets: 4DFR is ranked at the $97.8$th percentile
while its own best candidate sits at the $33.0$th, 1N46 at the $92.1$st against
$14.6$th, and 1YGC at the $95.8$th against $74.7$th.

\begin{table}[htbp]
\centering
\caption{Minimum Euclidean distance between the true ligand centroid and the top-ranked
candidate window ($k=3$) across the $29$ evaluated protein--ligand pairs.
\emph{Failed (Threshold)} indicates that no exact bitstring match cleared the $1/N$
amplitude threshold.}
\label{tab:distances}
\begin{tabular}{llll}
\toprule
\textbf{Target} & \textbf{Top-1 distance (\AA)} & \textbf{Target} & \textbf{Top-1 distance (\AA)} \\
\midrule
\texttt{1GM8} & $9.74$  & \texttt{1R9O} & $7.76$ \\
\texttt{1GPK} & $7.40$  & \texttt{1S19} & $8.64$ \\
\texttt{1HNN} & \emph{Failed (Threshold)} & \texttt{1STP} & \emph{Failed (Threshold)} \\
\texttt{1HQ2} & $6.93$  & \texttt{1T46} & $9.00$ \\
\texttt{1HSG} & $6.48$  & \texttt{1TZ8} & \emph{Failed (Threshold)} \\
\texttt{1IG3} & $4.26$  & \texttt{1U1C} & $6.52$ \\
\texttt{1K3U} & $7.62$  & \texttt{1XM6} & \emph{Failed (Threshold)} \\
\texttt{1L2S} & $6.49$  & \texttt{1Y6B} & $7.75$ \\
\texttt{1MMV} & $9.21$  & \texttt{1YGC} & $22.18$ \\
\texttt{1N2V} & \emph{Failed (Threshold)} & \texttt{1YV3} & $5.10$ \\
\texttt{1N46} & $12.08$ & \texttt{1YWR} & $8.22$ \\
\texttt{1OPK} & $11.93$ & \texttt{2BSM} & $11.20$ \\
\texttt{1OYT} & $5.00$  & \texttt{3PTB} & $6.81$ \\
\texttt{1Q1G} & $4.50$  & \texttt{4DFR} & $14.23$ \\
\texttt{1R58} & \emph{Failed (Threshold)} & & \\
\bottomrule
\end{tabular}
\end{table}

\textbf{Is the criterion reachable?} A distance of $7.75$~\AA{} is only meaningful once
it is known how close a correct answer could have come. The ligand centroid lies in the
middle of the cavity, in empty space, whereas pharmacophoric sites lie on the atoms that
line it, so there is a floor below which no window can fall. If that floor were above
$5$~\AA{} the criterion of Table~\ref{tab:distances} would be unreachable by
construction and the benchmark would be measuring pocket geometry rather than the
descriptor. We therefore computed, for every pocket, the distance to the ligand centroid
of the \emph{closest window that exists at all} --- the ceiling on any possible search
--- and of the \emph{median window}, which is the chance level expressed in \AA{}
rather than as a percentile. Both quantities range over all unit-stride windows, the same
population against which Table~\ref{tab:percentile} normalises, and neither involves the
oracle, the acceptance threshold or the ranking; they are properties of the
representation alone, and are defined even for the two pockets whose parsed ligand yields
no pharmacophoric sites. Table~\ref{tab:ceiling} reports them.

\begin{table}[htbp]
\centering
\caption{Reachability of the $5$~\AA{} criterion at $k=3$, over all unit-stride windows
of the $29$ evaluated pockets. \emph{Best window} is the closest window that exists in
the representation; \emph{median window} is the chance level in \AA{}. Medians
are taken across pockets.}
\label{tab:ceiling}
\begin{tabular}{lcc}
\toprule
\textbf{Window} & \textbf{Median} & \textbf{Range across pockets} \\
\midrule
Best available in the representation (ceiling) & $3.76$~\AA{} & $1.21$--$5.67$ \\
Median window of the pocket (chance)           & $7.88$~\AA{} & $6.39$--$10.99$ \\
Top-1 returned by the search (Table~\ref{tab:distances}) & $7.75$~\AA{} & $4.26$--$22.18$ \\
\midrule
\multicolumn{2}{l}{Pockets where a window within $5$~\AA{} exists} & $25/29$ \\
\multicolumn{2}{l}{Pockets where the search returns one, of those reachable} & $2/19$ \\
\bottomrule
\end{tabular}
\end{table}

The criterion is therefore fair: in $25$ of the $29$ pockets a window within $5$~\AA{} of
the ligand centroid is physically present in the representation, and the median target
sits at $3.76$~\AA{}. The four exceptions --- 1OPK, 1S19, 1Y6B and 1YWR, whose ceilings
lie between $5.01$ and $5.67$~\AA{} --- are large pockets in which no window of three
sites reaches the centre of the cavity, and they are the only cases where a failure at
$5$~\AA{} carries no information about the descriptor. All four returned a candidate and
are therefore counted among the $23$ pockets of Table~\ref{tab:distances}; excluding them
leaves $19$ pockets on which the criterion could in principle have been met, and the
success rate is properly read as $2/19$ rather than $2/23$.

The table also restates the result of Table~\ref{tab:percentile} in absolute units,
which is how it is most easily read: the target was at $3.76$~\AA{}, drawing a window at
random gives $7.88$~\AA{}, and the search returns $7.75$~\AA{}. The measured gap between
chance and the search is $0.13$~\AA{} on the useful side, against the $4.12$~\AA{} that
separates chance from the best answer available. The search recovers essentially none of
the distance it had room to recover, and it does so on pockets where that room
demonstrably exists.

% ───────────────────────────────────────────────────────────────────
\subsection{Neither the window size nor the matching regime recovers the signal}
\label{sec:res-sweep}

Because the matching set is determined classically, its quality can be measured
directly, independently of any ranking applied afterwards. Table~\ref{tab:ksweep}
reports the sweep over the window size $k$ in the exact regime.

\begin{table}[htbp]
\centering
\caption{Effect of the sliding-window size on exact-bitstring-match coverage and on the
spatial percentile of the matching set.}
\label{tab:ksweep}
\begin{tabular}{ccccc}
\toprule
\textbf{Max sites ($k$)} & \textbf{Coverage} & \textbf{Median match percentile} & \textbf{Best of set} & \textbf{Baseline} \\
\midrule
$3$ & $23/29$ & $48.2\%$ & $16.8\%$ & $18.4\%$ \\
$4$ & $12/29$ & $49.7\%$ & $11.7\%$ & $23.9\%$ \\
$5$ & $7/29$  & $31.1\%$ & $9.8\%$  & $39.9\%$ \\
\bottomrule
\end{tabular}
\end{table}

At $k=3$ the matching set sits at the $48.2$nd percentile --- exactly chance. Since a
ranking function reorders a set without adding information to it, this bounds what any
ranking criterion whatsoever could achieve at the default configuration. Raising $k$ to
$5$ does produce a measurable signal ($9.8\%$ against a $39.9\%$ expected baseline),
but coverage collapses from $23$ to $7$ pockets. The cause is informational rather than
chemical: the bitstring takes $4^k$ values, so at $k=3$ its $64$ codes are spread over
the $56$--$616$ windows of a typical pocket and collisions are frequent, whereas at
$k=5$ the $1024$ codes make an exact match rare --- rarer, in most pockets, than the
ligand's own code occurring at all.

Replacing exact equality with a distance in the second-encoding amplitude space
decouples the two effects. At $k=5$ this restores coverage from $7/29$ to $27/29$
pockets --- a genuine and worthwhile gain --- but leaves accuracy unchanged: the top-1
candidate lands at the $45.8$th percentile (median $42.1$th, fidelity-ranked $45.0$th),
against $50\%$ for chance, and the best of three candidates at the $28.1$th against
$25\%$ expected from three random draws. Coverage and accuracy are therefore governed by
different bottlenecks, and only the first is a property of the matching criterion.

% ───────────────────────────────────────────────────────────────────
\subsection{The sliding windows are spatially well formed}
\label{sec:res-geometry}

A natural suspicion is that $k$ consecutive sites of the flat chain are spatially
scattered, so that a window would not correspond to any physical patch of the pocket
surface and the comparison against a ligand centroid would be meaningless. Measured over
all windows of all pockets at $k=5$, this is not the case: the median window span is
$7.60$~\AA{} with a median radius of gyration of $3.06$~\AA{}, and only $26.1\%$ of
windows exceed $10$~\AA{}. Windows are compact, local surface patches of a size
comparable to a small ligand. The chain-assembly and topological-ordering stages of
Sections~\ref{sec:topological-order}--\ref{sec:chain-assembly} therefore do preserve spatial coherence,
as they were designed to.

% ───────────────────────────────────────────────────────────────────
\subsection{The descriptor carries no information about ligand position}
\label{sec:res-spearman}

The preceding subsections bound the performance of the search without explaining it. The
decisive measurement removes the search entirely and asks whether the descriptor and the
ground truth are related at all. For every window of every pocket we computed its
distance from the ligand in second-encoding amplitude space and its true 3D distance
from the ligand centroid, and correlated the two by Spearman's rank coefficient. At the
default window size $k=3$, over all 29 benchmark pockets:

\begin{center}
\begin{tabular}{lc}
\toprule
Mean $\rho$ across pockets      & $-0.008$ \\
Median $\rho$                   & $+0.013$ \\
Pockets with $|\rho| > 0.3$     & $2/29$ \\
Pockets with the correct sign   & $16/29$ \\
\bottomrule
\end{tabular}
\end{center}

Widening the window does not change the picture: at $k=5$ the mean is $+0.020$, the
median $+0.069$, and the sign is correct in 17 of 29 pockets.

The correlation is absent rather than weak, and the sign is correct in barely more than
half the pockets, as expected from noise. Two windows that encode similarly are no more likely to be
spatially close to each other, or to the ligand, than two windows chosen at random.
Since this quantity is measured over all windows and involves neither the oracle, nor
the acceptance threshold, nor the ranking, it is a property of the representation alone.

\todo[inline, color=yellow!40]{Da completare: sottosezione con i tre controlli
(degenerazione del codice, AUC contro la ground truth, controllo positivo con le
distanze fra siti) e il controllo di robustezza escludendo le \texttt{PosIonizable} di
backbone. Vedi \texttt{scripts/diagnose\_descriptor.py}.}
